% Source: https://artofproblemsolving.com/community/c3h1710142_junior_regional_olympiad__fbh_2017_grade_7_problem_5
% Stable identity: aops-topic-1710142-q01
Suppose $p$ is rational; clearly $p>0$. Then $\sqrt2-p\sqrt3=p\sqrt b-\sqrt a$. Squaring both sides,
\[
2+3p^2-2p\sqrt6=p^2b+a-2p\sqrt{ab},
\]
so $2p\left(\sqrt6-\sqrt{ab}\right)=2+3p^2-p^2b-a$ is rational, and therefore $r=\sqrt6-\sqrt{ab}$ is rational.

\textbf{$ab=6$.} From $\sqrt{ab}=\sqrt6-r$, squaring gives $ab=6-2r\sqrt6+r^2$. If $r\ne0$ this would make $\sqrt6$ rational, which it is not. So $r=0$, i.e.\ $ab=6$, and then also $2+3p^2=p^2b+a$.

\textbf{Check the four cases.}
\begin{itemize}
\item $(a,b)=(1,6)$: $p=\dfrac{\sqrt2+1}{\sqrt3+\sqrt6}=\dfrac{\sqrt2+1}{\sqrt3\,(1+\sqrt2)}=\dfrac{1}{\sqrt3}$, irrational.
\item $(a,b)=(2,3)$: $p=\dfrac{2\sqrt2}{2\sqrt3}=\sqrt{\dfrac{2}{3}}$, irrational.
\item $(a,b)=(3,2)$: $p=\dfrac{\sqrt2+\sqrt3}{\sqrt3+\sqrt2}=1$, rational.
\item $(a,b)=(6,1)$: $p=\dfrac{\sqrt2+\sqrt6}{\sqrt3+1}=\dfrac{\sqrt2\,(1+\sqrt3)}{\sqrt3+1}=\sqrt2$, irrational.
\end{itemize}
Hence the only solution is $\boxed{a=3,\ b=2}$ (with $p=1$).
